\SixthArticle{22}
Although the function $\Pi(k,z)$ may seem more general than $\Pi z$, it will from now on be unnecessary for us, since it is readily reduced to the latter. Combining equations 38, 45, 46 gives
\[
\begin{array}{@{}ll@{}}
[47]&\displaystyle\Pi(k,z)=\frac{k^z\Pi k.\Pi z}{\Pi(k+z)}
\end{array}
\]
Moreover, the connection of these functions with those that the celebrated \textsc{Kramp} called \emph{numerical\SourceLeaf{156} factorials} is evident in itself. Namely, the numerical factorial that this author denotes by
$a^{bIc}$
is, in our notation,


\[
=\frac{c^b b^{\frac{a}{c}-1}\Pi b}{\Pi(b,\frac{a}{c}-1)}
=\frac{c^b\Pi(\frac{a}{c}+b-1)}{\Pi(\frac{a}{c}-1)}
\]
But it seems more advisable to introduce a function of \emph{one} variable into analysis than a function of three variables, especially since the latter can be reduced to the former.

\SixthArticle{23}
The continuity of the function $\Pi z$ is interrupted whenever its value becomes infinitely large, that is, for negative integer values of $z$. Thus it will be positive from $z=-1$ to $z=\infty$, and since $\Pi z$ takes an infinitely large value at both limits, there will be a minimum value between them, which we have found to be $=0{,}8856024$, corresponding to the value $z=0{,}4616321$. Between the limits $z=-1$ and $z=-2$ the value of the function $\Pi z$ becomes negative, between $z=-2$ and $z=-3$ positive again, and so on, as follows at once from equation 44. Further, if all the values of the function $\Pi z$ between any two limits differing by unity, for example from $z=0$ to $z=1$, may be regarded as known, the value of the function for any other real value of $z$ can evidently be deduced from them with the aid of equation 45. For this purpose we have constructed a \emph{table}, appended at the end of this section, giving to twenty figures the Briggsian logarithms of the function $\Pi z$ for $z=0$ to $z=1$ at steps of one hundredth, calculated with the greatest care. It must nevertheless be noted that the last, twentieth figure may sometimes be in error by one or two units.

